El reporte debe subrayar varias limitaciones metodol\'ogicas para evitar una lectura excesiva de los resultados:

\begin{itemize}
\item No hay una funci\'on de completitud instrumental completa por estrella, m\'etodo o campa\~na observacional.
\item Los candidatos generados son interpolaciones orbitales plausibles, no detecciones ni soluciones din\'amicas confirmadas.
\item Las masas y radios candidatos se estiman por interpolaci\'on local o por an\'alogos globales; no son medidas directas.
\item Un gap amplio puede reflejar historia de formaci\'on, inestabilidad orbital, migraci\'on o arquitectura real del sistema, no solo incompletitud observacional.
\item TOI, ATI y sombra provienen de un Mapper global y no prueban causalidad f\'isica dentro de un sistema particular.
\item El proxy RV no equivale a una amplitud de velocidad radial completa con inclinaci\'on, excentricidad y ruido estelar.
\item El proxy de tr\'ansito no reemplaza una modelaci\'on fotom\'etrica real ni incorpora ventanas temporales, cadencia o ruido instrumental.
\item El m\'odulo no incorpora de forma expl\'icita inclinaci\'on, excentricidad, actividad estelar, sensibilidad instrumental detallada ni estabilidad N-body.
\end{itemize}

Estas limitaciones no invalidan el valor del m\'odulo, pero s\'i delimitan su alcance correcto: una herramienta de se\~nal topol\'ogica exploratoria y priorizaci\'on observacional, no un detector definitivo de planetas ausentes.
